\section*{Chapitre \ref{ch:iv1:machine-learning} --- Apprentissage automatique}
\markboth{\protect\omTocMark{}{Réponses~: Apprentissage automatique}}{\protect\omTocMark{}{Réponses~: Apprentissage automatique}}

\begin{solution}{iq:iv1:machine-learning:1}
Exactitude $0.93$, précision $0.40$, rappel $0.80$, F1 $0.53$. Un modèle qui ne signale jamais rien atteint une exactitude de $0.95$, valeur supérieure.

Les effectifs sont 40 vrais positifs, 60 faux positifs, 10 faux négatifs, et donc 890 vrais négatifs. L'exactitude vaut $930/1000$. La précision\index{précision et rappel} vaut $40/100$, la part des signalements qui sont justes. Le rappel vaut $40/50$, la part des événements réels qui sont détectés. Le F1\index{score F1} est leur moyenne harmonique, $2(0.4)(0.8)/1.2 = 0.533$.

Le modèle trivial a raison sur les 950 jours sans événement, et bat donc le vrai modèle en exactitude tout en étant inutile. Lorsque les classes sont rares, l'exactitude mesure le taux de base.

\iqwrong{Présenter une exactitude de $0.93$ comme un bon résultat. Ou additionner 40 et 10 pour obtenir les jours signalés, ce qui confond les deux dénominateurs.}

\iqfollow{Écrivez le F1 sans faire référence aux vrais négatifs.}{$\mathrm{TP}/(\mathrm{TP} + \frac12(\mathrm{FP} + \mathrm{FN})) = 40/75$. Les 890 vrais négatifs n'ont aucun effet, et c'est précisément l'intérêt de la mesure.}

\iqwhy{C'est le contrôle le plus rapide pour vérifier qu'un candidat lit les bons dénominateurs, et l'exemple du modèle inutile montre si l'exactitude sera remise en question.}
\end{solution}

\begin{solution}{iq:iv1:machine-learning:2}
Une exactitude d'environ 59~\%, alors que chaque modèle vaut en réalité 50~\%. Un ensemble de test ne s'utilise qu'une seule fois, une fois toute la sélection terminée.

Le score de chaque modèle est à peu près normal, de moyenne 50~\% et d'écart-type 5~\%. Le maximum de 20 normales indépendantes se situe environ $1.87$ écarts-types au-dessus de la moyenne, si bien que le meilleur attendu vaut environ $59.3~\%$, et un calcul binomial exact comme une simulation donnent la même valeur.

La bonne procédure~: choisir sur un ensemble de validation distinct, ou par validation croisée\index{validation croisée}, puis évaluer une seule fois le modèle retenu sur l'ensemble de test resté intact. Ce score est sans biais pour le modèle retenu, et ici il reviendrait proche de 50~\%. Si plusieurs modèles doivent être présentés, il faut tenir compte du nombre de modèles essayés.

\iqwrong{Présenter le score de test du meilleur modèle, ou utiliser l'ensemble de test pour décider de la suite à essayer. Chaque coup d'œil transforme l'ensemble de test en ensemble de validation.}

\iqfollow{Pendant combien de jours un ensemble de test permet-il de garder le meilleur de 20 modèles nuls sous 52~\% en moyenne~?}{Environ 2200 jours~: $1.87 \times 50/\sqrt{n} = 2$ donne $n \approx 2180$.}

\iqwhy{Tout le sujet des recherches sur le surapprentissage tient dans cet exemple. Les recruteurs quant écoutent pour savoir si le candidat compte le nombre d'essais sans qu'on le lui demande.}
\end{solution}

\begin{solution}{iq:iv1:machine-learning:3}
La régression ridge, les $k$ plus proches voisins et la régression sur composantes principales changent. Les moindres carrés ordinaires, l'arbre de décision et la forêt aléatoire ne changent pas, hormis les éventuelles graines aléatoires.

Les moindres carrés multiplient le coefficient par $1/100$, et les prédictions ne changent pas. La régression ridge\index{régression ridge} pénalise la taille des coefficients~: un coefficient plus petit sur la variable remise à l'échelle coûte moins cher, et l'ajustement se déplace. Les plus proches voisins utilisent des distances, si bien que la variable exprimée en centimes les domine. Les composantes principales maximisent la variance, si bien que cette même variable prend la première composante.

Un arbre de décision ne découpe que selon l'ordre, et multiplier par 100 conserve cet ordre~; une forêt aléatoire est un ensemble d'arbres. Des essais sur données simulées confirment chaque cas.

\iqwrong{Dire que la régression ridge n'est pas touchée parce que les moindres carrés ne le sont pas. Ou affirmer que tous les modèles devraient bien se comporter après l'ajustement. La règle est~: tout ce qui utilise une distance ou une pénalité sur la taille exige une échelle commune.}

\iqfollow{Que devez-vous faire de tels modèles, et à quel moment~?}{Standardiser chaque variable, en utilisant uniquement les données d'entraînement. Utiliser toutes les données introduit une fuite d'information depuis la période de test.}

\iqwhy{Une question de préparation élémentaire qui piège ceux qui traitent tous les modèles de la même façon, et c'est une vraie source de bogues.}
\end{solution}

\begin{solution}{iq:iv1:machine-learning:4}
La constante $0.1$ obtient une perte logarithmique de $0.325$, la constante $0.5$ obtient $0.693$, et la ligne unique coûte $\ln 1000 = 6.91$.

La perte logarithmique\index{perte logarithmique} est la moyenne de $-\ln(\text{probabilité prédite de ce qui s'est produit})$. Pour la constante $0.1$~: $-0.1\ln 0.1 - 0.9\ln 0.9 = 0.325$. Pour $0.5$~: $-\ln 0.5 = 0.693$, pour chaque ligne, quel que soit le taux de base.

Une prédiction de $0.999$ sur une étiquette 0 accorde la probabilité $0.001$ à ce qui s'est produit, et coûte $6.91$. C'est plus de vingt fois la perte de référence, et une ligne de ce type sur 100 ajoute $0.069$, soit un cinquième de la perte de référence. La métrique sanctionne donc sévèrement la surconfiance. Un modèle qui renvoie des scores bruts proches de 0 et de 1 a besoin d'un écrêtage ou d'une calibration.

\iqwrong{Comparer la perte logarithmique d'un modèle à $0.693$ sur un problème déséquilibré. La référence juste est la constante égale au taux de base, ici $0.325$.}

\iqfollow{Quelle constante minimise la perte logarithmique~?}{Le taux de base. En annulant la dérivée de $-b\ln p - (1 - b)\ln(1 - p)$, on obtient $p = b$, ce qui explique pourquoi la perte logarithmique récompense la calibration.}

\iqwhy{Lire une perte demande une référence. Un candidat qui donne immédiatement le chiffre du taux de base a utilisé cette métrique sur des données réelles.}
\end{solution}

\begin{solution}{iq:iv1:machine-learning:5}
La cote double, puisque $e^{0.7} = 2.01$. Une probabilité de $0.2$ devient $0.33$, et une probabilité de $0.5$ devient $0.67$.

Le modèle est linéaire en log-cote\index{log-cote}, donc une hausse d'une unité ajoute $0.7$ à la log-cote et multiplie la cote par $e^{0.7}$.

À partir de $0.2$, la cote vaut $0.25$, qui double pour donner $0.5$, et donc $p = 0.5/1.5 = 1/3$. À partir de $0.5$, la cote vaut 1, qui double pour donner 2, et donc $p = 2/3$. Le même coefficient ajoute $0.13$ dans un cas et $0.17$ dans l'autre. L'effet sur la probabilité dépend du point de départ.

\iqwrong{Dire que la probabilité augmente de $0.7$, ou qu'elle double. Une probabilité ne peut pas le faire à partir de $0.6$, alors que la cote le peut.}

\iqfollow{Donnez une règle rapide pour l'effet au voisinage de $p = \frac12$.}{La pente de la courbe logistique y vaut $\frac14$, donc la variation vaut environ le coefficient divisé par 4. Ici cela donne $0.175$, proche de la valeur exacte $0.17$.}

\iqwhy{On lit les coefficients de travers plus souvent qu'on ne les calcule de travers. Distinguer la cote de la probabilité est précisément ce que teste la question.}
\end{solution}

\begin{solution}{iq:iv1:machine-learning:6}
Les ensembles d'entraînement comptent 800 lignes. Une ligne est validée une fois et contribue à l'entraînement de quatre modèles. La recherche nécessite 5000 ajustements.

Chaque pli met de côté un cinquième, soit 200 lignes, et l'entraînement porte sur les 800 autres. Sur les cinq plis, une ligne figure exactement une fois dans la partie mise de côté et quatre fois dans la partie d'entraînement.

La grille compte $10^3 = 1000$ points, et chacun demande 5 ajustements, soit 5000 au total\index{recherche d'hyperparamètres}. Une fois le point choisi, un ajustement supplémentaire sur les 1000 lignes donne le modèle final.

Le score de validation croisée du point retenu est biaisé à la hausse, puisqu'il est le meilleur sur 1000. Il faut un ensemble de test distinct, ou une boucle externe, pour obtenir un chiffre honnête.

\iqwrong{Présenter le score du meilleur point de la grille comme performance attendue, ou multiplier 3 par 10 pour la taille de la grille. Il s'agit de 10 à la puissance 3.}

\iqfollow{Combien coûte une recherche imbriquée à 5 plis externes~?}{$25\,000$ ajustements, cinq fois plus. C'est le coût d'une estimation sans biais de toute la procédure.}

\iqwhy{Des décomptes simples révèlent qui a réellement mené une telle recherche, et qui a réfléchi à ce que signifie son meilleur score.}
\end{solution}

\begin{solution}{iq:iv1:machine-learning:7}
Environ 15~\% des signalements sont réels, 48~\% au taux de faux positifs plus faible, et l'aire sous la courbe ROC ne change pas avec la prévalence.

Sur 1000 transactions~: 10 sont mauvaises, dont 9 sont détectées, et 990 sont bonnes, dont $49.5$ sont signalées. La précision vaut donc $9/58.5 = 0.154$. Avec un taux de faux positifs de 1~\%, les signalements valent $9 + 9.9$, ce qui donne $0.476$. Une petite amélioration du taux de faux positifs a un effet important, parce que les bonnes transactions sont bien plus nombreuses que les mauvaises.

La courbe ROC représente le taux de vrais positifs en fonction du taux de faux positifs\index{courbe ROC}. Chacun se calcule au sein d'une seule classe, si bien que la composition des classes n'intervient pas. Un calcul pour les deux prévalences donne la même aire.

\iqwrong{Dire qu'un détecteur avec des taux de 90~\% et de 95~\% a raison la plupart du temps lorsqu'il signale. Ou présenter seulement l'aire ROC pour un problème d'événements rares, où elle masque la charge de fausses alarmes.}

\iqfollow{Quelle est la part des signalements réels à une prévalence de 10~\%~?}{$0.09/(0.09 + 0.045) = 0.667$, avec le même détecteur. La précision dépend bien de la prévalence, si bien qu'un chiffre mesuré sur un ensemble de test équilibré ne se transpose pas en production.}

\iqwhy{Elle distingue les métriques qui dépendent de la prévalence de celles qui n'en dépendent pas, et elle revient dès qu'un modèle passe à un flux de données à la composition différente.}
\end{solution}

\begin{solution}{iq:iv1:machine-learning:8}
Il y a 3457 paramètres, soit environ $2.8$ par ligne, si bien qu'avec des données peu informatives le modèle mémorisera, à moins d'être fortement contraint.

Chaque couche contient le produit du nombre d'entrées par le nombre d'unités pour les poids, plus un biais par unité. Première couche~: $(20 + 1) \times 64 = 1344$. Deuxième couche~: $(64 + 1) \times 32 = 2080$. Sortie~: $(32 + 1) \times 1 = 33$. La somme vaut 3457, et le décompte donné par une bibliothèque d'apprentissage profond le confirme.

Avec 1250 lignes, le modèle compte davantage de paramètres que de lignes, et les rendements ne sont pour l'essentiel que du bruit, si bien qu'il peut ajuster l'ensemble d'entraînement sans aucun signal réel. La décroissance des poids, le dropout, l'arrêt précoce et une architecture plus petite sont les réponses habituelles. Une validation walk-forward honnête tranche. Un réseau de neurones très petit, ou un modèle linéaire, constitue la référence à battre.

\iqwrong{Oublier les biais, ce qui donne 3360. Ou conclure que le réseau doit forcément échouer~: des réseaux surparamétrés peuvent généraliser, mais sur des rendements bruités, la taille du modèle est un signal d'alerte qui demande à être testé.}

\iqfollow{Que tenteriez-vous en premier avec si peu de données~?}{Un modèle linéaire régularisé, comme référence. Si un réseau ne parvient pas à le battre hors échantillon, les paramètres supplémentaires ne rapportent rien.}

\iqwhy{Compter les paramètres est un exercice mécanique, et ce qui retient l'attention du recruteur, c'est la comparaison avec le nombre de lignes.}
\end{solution}

\begin{solution}{iq:iv1:machine-learning:9}
Le modèle A surapprend, avec une variance élevée. Le modèle B sous-apprend, avec un biais élevé. Davantage de données aident A et non B.

Le modèle A présente un écart de 29~points entre l'entraînement et la validation. Il mémorise. Davantage de données réduisent l'écart\index{courbes d'apprentissage}, comme le font une régularisation plus forte, moins de variables et un modèle plus simple.

Le modèle B présente presque aucun écart, et les deux courbes restent très en dessous de la cible de 85~\%. Il lui manque de la capacité ou de l'information. Ajouter des lignes déplace très peu les deux courbes. Il lui faut un modèle plus souple, ou de meilleures variables. Si la cible est hors d'atteinte pour tout modèle, le bruit des étiquettes est peut-être la vraie limite~: c'est donc ce qu'il faut vérifier d'abord.

Une simulation montre les deux cas. Un classifieur au plus proche voisin sur une sinusoïde bruitée présente un large écart qui se réduit avec davantage de lignes. Un classifieur linéaire sur la même sinusoïde reste bas sur les deux ensembles et ne s'améliore pas.

\iqwrong{Dire que davantage de données aide toujours. Cela ne fait rien pour un modèle déjà trop simple, puisque les courbes de B se sont déjà rejointes.}

\iqfollow{Que vous suggérerait un score de validation supérieur au score d'entraînement~?}{Une régularisation active seulement à l'entraînement, comme le dropout, ou une période de validation plus facile. Si aucune de ces causes ne s'applique, soupçonner un découpage qui fuit.}

\iqwhy{Le diagnostic décide où va le temps. Un candidat qui sait dire quelle correction serait du temps perdu fait gagner des semaines de travail.}
\end{solution}

\begin{solution}{iq:iv1:machine-learning:10}
Le biais vaut $0.1(k + 1)$ et la variance $1/k$, si bien que le meilleur $k$ vaut 3, avec une erreur quadratique de $0.493$. Doubler la densité porte le meilleur $k$ à 6.

La moyenne des $k$ voisins s'écarte en espérance de $0.2(1 + \dots + k)/k = 0.1(k + 1)$ de la vraie valeur\index{biais et variance}. Sa variance vaut $1/k$. L'erreur quadratique est donc
\[
\frac{1}{k} + 0.01(k + 1)^2.
\]
Les valeurs pour $k = 1, 2, 3, 4$ sont $1.04$, $0.59$, $0.493$ et $0.5$, si bien que le meilleur est $k = 3$. Une simulation de l'estimateur à trois voisins donne le même $0.493$.

Avec un espacement de $0.1$, l'erreur vaut $1/k + 0.0025(k + 1)^2$, avec un minimum de $0.289$ en $k = 6$. Davantage de données permet un voisinage plus grand, ce qui abaisse la variance tant que les distances restent faibles. Le biais et la variance diminuent tous deux.

Chaque nouvelle observation porte encore un bruit de variance 1, que nul choix de $k$ ne supprime.

\iqwrong{Ajouter le biais à la variance sans le mettre au carré. Les termes qui s'additionnent sont la variance et le carré du biais.}

\iqfollow{Si les voisins se répartissent également de part et d'autre, quel est le meilleur $k$~?}{Aussi grand que les données le permettent, dans ce modèle. Une droite donne un biais nul sous une moyenne symétrique, si bien que seule la variance subsiste. C'est la courbure qui l'empêcherait.}

\iqwhy{Une courbe biais-variance explicite est plus convaincante que la figure habituelle, et elle montre comment la complexité optimale se déplace avec les données.}
\end{solution}

\begin{solution}{iq:iv1:machine-learning:11}
Erreur d'entraînement $0.8$, erreur sur un nouveau bruit $1.2$. L'écart de $0.4$ vaut $2p\sigma^2/n$.

L'ajustement projette le bruit sur l'espace engendré par les régresseurs. L'ajustement absorbe $p$ des $n$ dimensions du bruit, si bien que la somme des carrés des résidus a pour espérance $\sigma^2(n - p)$\index{optimisme}, et que l'erreur quadratique moyenne vaut $1 - 20/100 = 0.8$.

Les valeurs ajustées portent un bruit de variance totale $p\sigma^2$. Sur un bruit neuf aux mêmes points, l'erreur espérée est ce bruit neuf auquel s'ajoute celui-ci~: $1 + 20/100 = 1.2$. Une simulation donne $0.80$ et $1.20$.

L'ajustement paraît donc meilleur que le niveau de bruit, et prédit moins bien d'autant. L'écart est linéaire en fonction du nombre de paramètres. C'est ce que le $C_p$ de Mallows et le critère d'Akaike ajoutent à l'erreur d'entraînement pour la corriger.

\iqwrong{Dire qu'un modèle correct ne surapprend pas. Avoir le bon modèle supprime le biais, non la variance liée à l'estimation de 20 nombres.}

\iqfollow{Que rajoutent à l'écart dix paramètres inutiles de plus~?}{$2 \times 10/100 = 0.2$. Un critère n'accepte une nouvelle variable que si elle abaisse l'erreur d'entraînement de plus que cette pénalité.}

\iqwhy{C'est le résultat exact le plus simple derrière la sélection de modèle. Il montre que l'erreur d'entraînement trompe même lorsque le modèle est le bon.}
\end{solution}

\begin{solution}{iq:iv1:machine-learning:12}
L'impureté de Gini baisse de $0.163$, et l'entropie de $0.256$ bits.

Fils gauche~: 30 et 10~; fils droit~: 10 et 50, avec les poids $0.4$ et $0.6$.

L'impureté de Gini\index{impureté de Gini} vaut $2p(1 - p)$. Nœud parent~: $2(0.4)(0.6) = 0.48$. Fils gauche~: $2(0.75)(0.25) = 0.375$. Fils droit~: $2(\frac16)(\frac56) = 0.278$. Moyenne pondérée~: $0.4(0.375) + 0.6(0.278) = 0.317$. La baisse vaut donc $0.163$.

L'entropie\index{entropie} vaut $-p\log_2 p - (1 - p)\log_2(1 - p)$. Nœud parent~: $0.971$. Fils gauche~: $0.811$. Fils droit~: $0.650$. Moyenne pondérée~: $0.715$. La baisse vaut $0.256$.

Les deux critères n'ont pas la même échelle, si bien que les nombres ne sont pas comparables, et en pratique ils classent la plupart des découpages de la même façon.

\iqwrong{Faire la moyenne des deux fils avec des poids égaux, ce qui donne un autre nombre. Les poids sont les parts de lignes.}

\iqfollow{Quelle est la signification de l'impureté de Gini~?}{La probabilité que deux lignes tirées au hasard avec remise dans le nœud portent des étiquettes différentes, ou l'erreur d'une règle qui devine au hasard selon les proportions du nœud.}

\iqwhy{Un contrôle simple~: savoir calculer ce qu'une bibliothèque calcule, la pondération étant l'erreur la plus fréquente.}
\end{solution}

\begin{solution}{iq:iv1:machine-learning:13}
L'erreur de classement donne $0.25$ pour les deux. L'impureté de Gini donne $0.375$ pour A et $0.333$ pour B, si bien que B est préféré, puisqu'il crée un nœud pur. La raison en est que Gini est strictement concave.

Sous la règle majoritaire~: dans A, chaque fils se trompe sur 100, soit $200/800$. Dans B, le premier fils se trompe sur 200 et le second sur 0, soit encore $200/800$.

Gini pour A~: $2(\frac34)(\frac14) = 0.375$ pour chaque fils. Pour B~: le premier fils a $\frac34$ des lignes, avec $2(\frac13)(\frac23) = 0.444$, et le fils pur vaut 0, donc $0.75 \times 0.444 = 0.333$.

Le taux d'erreur de classement est linéaire par morceaux en $p$\index{critère de découpage}. Sauf si la majorité change d'un fils à l'autre, les deux fils donnent la même erreur pondérée que le parent, si bien que le découpage ne montre aucun gain. Une fonction strictement concave baisse dès qu'un découpage sépare les classes, même un peu, et elle récompense la pureté. Le nœud pur a de la valeur, car les découpages suivants n'ont plus qu'à traiter le reste.

\iqwrong{Dire que A est meilleur parce qu'il est symétrique, ou conclure que les critères sont toujours en désaccord. Ils s'accordent en général, mais les égalités et quasi-égalités, où ils peuvent diverger, sont fréquentes au début d'un arbre. Les fils $(210, 400)$ et $(190, 0)$, par exemple, ont une erreur plus élevée que A, $0.2625$, et une impureté de Gini plus faible, $0.344$.}

\iqfollow{Dans quel cas un découpage donne-t-il un gain nul, pour tout critère concave~?}{Lorsque les deux fils ont les proportions du parent. Dans ce cas, on n'a rien appris sur l'étiquette.}

\iqwhy{Cela explique pourquoi un arbre n'optimise pas simplement la métrique utilisée au final, et c'est un principe général utile.}
\end{solution}

\begin{solution}{iq:iv1:machine-learning:14}
La variance vaut $\rho\sigma^2 + (1 - \rho)\sigma^2/B$, soit $0.307$ pour $B = 100$, et $0.3$ à la limite. Le compromis est favorable, avec $0.131$ pour le même $B$.

Développons la variance de la moyenne\index{bagging}~:
\[
\frac{1}{B^2}\left(B\sigma^2 + B(B - 1)\rho\sigma^2\right) = \rho\sigma^2 + \frac{(1 - \rho)\sigma^2}{B}.
\]
Avec les valeurs énoncées~: $0.3 + 0.7/100 = 0.307$. On vérifie les limites~: $\rho = 1$ donne $\sigma^2$, comme un seul modèle, et $\rho = 0$ donne $\sigma^2/B$. Une simulation de modèles corrélés donne $0.307$.

Ajouter des modèles cesse donc d'aider à partir de la part commune $\rho\sigma^2$. Pour descendre plus bas, il faut rendre les modèles moins semblables, et c'est le rôle du sous-échantillonnage des variables dans une forêt aléatoire\index{forêt aléatoire}. Avec $\rho = 0.1$ et $\sigma^2 = 1.2$~: $0.12 + 1.08/100 = 0.131$. Les modèles individuels sont moins bons, alors que la moyenne est plus de deux fois meilleure.

\iqwrong{Supposer que la variance tend vers zéro avec assez d'arbres. Cela n'est vrai que pour des modèles indépendants, ce que ne sont pas les arbres construits sur les mêmes données.}

\iqfollow{Avec $\rho = 0.3$, combien de modèles suffisent pour s'approcher de la limite à 1~\% près~?}{$0.7/B \le 0.003$ donne environ 234. Au-delà, en ajouter davantage est du gaspillage.}

\iqwhy{La formule explique pourquoi les forêts fonctionnent et pourquoi les ensembles d'alphas semblables déçoivent. Les quants rencontrent la même formule lorsqu'ils combinent des signaux corrélés.}
\end{solution}

\begin{solution}{iq:iv1:machine-learning:15}
La part ajustée après $m$ arbres vaut $1 - 0.9^m$, qui dépasse 90~\% pour $m = 22$ et 99~\% pour $m = 44$. Le rétrécissement moyenne le bruit.

Pour $L = \frac12(y - F)^2$, la dérivée par rapport à $F$ vaut $-(y - F)$. Le gradient négatif est donc le résidu, et ajuster un arbre aux résidus revient à faire un pas de gradient dans l'espace des fonctions\index{gradient boosting}.

Si chaque arbre ajuste exactement le résidu et est ajouté avec le poids $\nu$, le résidu est multiplié par $1 - \nu = 0.9$ à chaque tour. Après $m$ tours, il reste $0.9^m$. Comme $0.9^{22} = 0.098$ et $0.9^{44} = 0.0097$, les deux seuils sont franchis à 22 et 44 tours.

Un arbre réel ajuste une estimation bruitée du résidu. Avec $\nu = 1$, tout le bruit du dernier arbre reste dans l'ajustement. Avec un $\nu$ petit, l'ajustement est une moyenne sur de nombreux arbres, chacun avec un bruit différent. Le prix à payer est un nombre d'arbres plus élevé.

\iqwrong{Croire que le rétrécissement modifie la cible finale. Il ne change que la vitesse à laquelle on l'atteint et la quantité de bruit qui l'accompagne.}

\iqfollow{Si l'estimation du résidu fournie par chaque arbre comporte un bruit de variance 1, indépendamment, à quelle variance de bruit l'ajustement se stabilise-t-il~?}{$\nu/(2 - \nu)$, soit $0.053$ pour $\nu = 0.1$ et 1 pour $\nu = 1$. Une simulation le confirme.}

\iqwhy{Un candidat doit savoir ce que fait le boosting au niveau d'une seule ligne. La formule explique le conseil habituel sur les petits taux d'apprentissage avec davantage d'arbres.}
\end{solution}

\begin{solution}{iq:iv1:machine-learning:16}
L'erreur de test de la forêt se stabilise et y reste, un peu au-dessus du bruit de 20~\%. L'erreur d'entraînement du boosting tend vers zéro, tandis que son erreur de test baisse puis remonte.

Une forêt moyenne des arbres. Ajouter un arbre ne modifie que l'estimation d'une moyenne, dont la variance décroît avec $B$, si bien que l'erreur tend vers une limite. Elle ne peut pas surapprendre par un excès d'arbres\index{forêt aléatoire}, bien que chaque arbre profond ait ajusté le bruit, et que l'erreur d'entraînement de la forêt soit elle aussi proche de zéro.

Le boosting est séquentiel\index{gradient boosting}. Chaque arbre vise ce qui reste inexpliqué, et une fois le motif épuisé, il ne reste plus que les étiquettes inversées. Le modèle a une capacité sans borne en nombre de tours, et il continue de les ajuster.

Un essai sur 400 lignes, avec un motif en damier et 20~\% d'étiquettes inversées, le montre. Le boosting avec un taux d'apprentissage de $0.5$ atteint une erreur d'entraînement inférieure à 1~\%, et son erreur de test finit au-dessus de 25~\%, soit plus de deux points au-dessus de son minimum. Pour la forêt, passer de 100 à 1000 arbres déplace l'erreur de test de moins d'un point. Donc, en boosting, le nombre d'arbres est un hyperparamètre à régler par arrêt précoce, alors que pour une forêt c'est une question de calcul.

\iqwrong{Dire que les deux surapprennent avec davantage d'arbres, ou qu'une forêt ne surapprend jamais. Elle ne surapprend pas par le nombre d'arbres, mais elle le peut par des arbres profonds sur des données bruitées.}

\iqfollow{Comment choisissez-vous le nombre de tours sur des données financières~?}{Par arrêt précoce sur un bloc de temps postérieur, séparé des données d'entraînement par un intervalle, et non sur un sous-ensemble aléatoire des lignes.}

\iqwhy{C'est le contraste central entre les deux familles. Le recruteur veut la raison, et pas seulement le fait.}
\end{solution}

\begin{solution}{iq:iv1:machine-learning:17}
$\alpha = 0.424$. Les lignes mal classées sont multipliées par $1.528$ et les lignes bien classées par $0.655$. Après normalisation, les lignes mal classées portent un poids total exactement égal à $\frac12$.

$\alpha = \frac12\ln(0.7/0.3) = 0.424$\index{AdaBoost}. Les facteurs sont $e^{\alpha} = \sqrt{7/3}$ et $e^{-\alpha} = \sqrt{3/7}$.

Les lignes mal classées portaient $0.3$, elles portent donc maintenant $0.3\sqrt{7/3} = \sqrt{0.21}$. Les lignes bien classées portaient $0.7$, elles portent maintenant $0.7\sqrt{3/7} = \sqrt{0.21}$. Ces deux quantités sont égales, donc chaque part vaut $\frac12$. Le facteur de normalisation est $2\sqrt{0.21} = 0.917$.

Cela signifie que le dernier apprenant a désormais une erreur pondérée égale à exactement $\frac12$ sur les nouveaux poids~: il ne vaut rien. Le suivant doit donc trouver quelque chose de nouveau. Il ne peut pas réussir en répétant le dernier. Cela vaut quelle que soit la valeur de $\varepsilon$, et c'est la raison de ce choix de $\alpha$.

\iqwrong{Dire que les lignes mal classées reçoivent «~plus de poids~» sans chiffre. Le $\frac12$ exact est le contenu même de l'algorithme.}

\iqfollow{Quels sont les votes pour $\varepsilon = 0.5$ et $\varepsilon = 0.6$~?}{Zéro et $\frac12\ln(2/3) = -0.20$. Un pile ou face ne reçoit aucun vote, et un apprenant pire qu'un pile ou face reçoit un vote négatif, de sorte que son opposé est utilisé.}

\iqwhy{Elle teste si la repondération se maîtrise comme un calcul, et elle mène à la borne sur l'erreur d'entraînement.}
\end{solution}

\begin{solution}{iq:iv1:machine-learning:18}
$\lambda = \sigma^2/\tau^2$. Le nombre effectif de paramètres vaut $\sum d_i^2/(d_i^2 + \lambda) = 1.5$.

La log-postérieure vaut, à une constante près,
\[
-\frac{1}{2\sigma^2}\|y - X\beta\|^2 - \frac{1}{2\tau^2}\|\beta\|^2.
\]
En multipliant par $2\sigma^2$, on obtient l'objectif de la régression ridge\index{régression ridge}, avec $\lambda = \sigma^2/\tau^2$. La postérieure est normale, donc son mode est sa moyenne, $(X^\top X + \lambda I)^{-1}X^\top y$. Un a priori plus serré ou des données plus bruitées donnent davantage de rétrécissement.

L'ajustement est $\hat{y} = Hy$ avec $H = X(X^\top X + \lambda I)^{-1}X^\top$. La définition naturelle des paramètres effectifs\index{degrés de liberté effectifs} est la trace de $H$, puisque pour les moindres carrés celle-ci vaut $p$. En fonction des valeurs singulières $d_i$ de $X$, elle vaut $\sum d_i^2/(d_i^2 + \lambda)$.

Avec les valeurs $100, 10, 1$ et $\lambda = 10$~: $100/110 + 10/20 + 1/11 = 0.909 + 0.5 + 0.091 = 1.5$. Une matrice ayant ces valeurs singulières donne la même trace. Pour $\lambda \to 0$ la somme vaut 3, soit le nombre complet de paramètres, et pour $\lambda \to \infty$ elle vaut 0. La direction forte garde sa liberté, et la direction faible est presque supprimée.

\iqwrong{Traiter $\lambda$ comme un simple bouton de réglage sans unité. Il dépend de l'échelle des variables et du bruit, ce que la forme bayésienne rend évident.}

\iqfollow{Quelle est la valeur de $\lambda$ pour un bruit de 1 et un écart-type a priori de $0.1$~?}{$\lambda = 1/0.01 = 100$. Pour les rendements, où les signaux sont faibles, de grandes valeurs de $\lambda$ sont normales et ne signalent pas une erreur.}

\iqwhy{Les deux points de vue sur la ridge, pénalité et a priori, sont souvent demandés ensemble. Les paramètres effectifs permettent au candidat de comparer un ajustement ridge à un ajustement non pénalisé.}
\end{solution}

\begin{solution}{iq:iv1:machine-learning:19}
Sans pénalité, $\beta$ part vers l'infini. Avec $\lambda = 0.1$, $\beta \approx 2.13$, et la probabilité en $x = 1$ vaut environ $0.89$.

La perte vaut $2\ln(1 + e^{-\beta})$, qui décroît indéfiniment quand $\beta$ augmente, et tend vers 0 sans minimiseur. Les données séparables\index{données séparables} n'ont pas d'estimateur du maximum de vraisemblance fini. En pratique le coefficient croît jusqu'à ce que le solveur s'arrête, et le logiciel peut avertir ou non.

Avec la pénalité, on annule la dérivée~:
\[
-\frac{2}{1 + e^{\beta}} + 0.1\beta = 0.
\]
Cette équation admet une racine unique en $\beta = 2.13$, trouvée par dichotomie et confirmée par un ajustement logiciel avec la même pénalité. La probabilité vaut $1/(1 + e^{-2.13}) = 0.894$.

La pénalité a donc ici un sens direct~: elle fixe le degré de confiance que le modèle peut avoir en l'absence de données contradictoires. Le résultat est une probabilité finie et raisonnable, obtenue à partir d'un jeu de données de deux lignes.

\iqwrong{Présenter l'ajustement sans pénalité comme un $\beta$ grand mais fini, ou comme une probabilité de 1. Dire aussi que les données séparables sont un bon signe. Avec de nombreuses variables et peu de lignes, c'est souvent simplement le signe d'un surapprentissage.}

\iqfollow{Que se passe-t-il pour $\lambda = 0.01$~?}{$\beta \approx 3.91$, avec une probabilité de $0.98$. Des pénalités plus faibles permettent davantage de confiance, mais lentement~: une division par dix de $\lambda$ ne double même pas $\beta$.}

\iqwhy{C'est une surprise numérique classique. Le candidat doit savoir pourquoi un ajustement logistique ne converge pas, et ce que la régularisation y apporte.}
\end{solution}

\begin{solution}{iq:iv1:machine-learning:20}
La vraie probabilité a posteriori vaut $0.23$ et le classifieur annonce $0.73$, si bien que la décision bascule.

Le Bayes naïf\index{Bayes naïf} multiplie les rapports de vraisemblance, en supposant l'indépendance conditionnellement à la classe. Avec un a priori de $\frac12$, la cote a posteriori est le produit des rapports.

La cote vraie utilise A une seule fois~: $3 \times 0.1 = 0.3$, donc la probabilité a posteriori vaut $0.3/1.3 = 0.23$. En dessous de $\frac12$, la décision est la classe 0.

Le classifieur compte A trois fois~: $3^3 \times 0.1 = 2.7$, donc la probabilité a posteriori vaut $2.7/3.7 = 0.73$. Au-dessus de $\frac12$, la décision est la classe 1. Un rapport de vraisemblance réel de 3 a été traité comme 27.

Avec A seul, le problème est masqué~: la vraie probabilité a posteriori vaut $0.75$, et celle du classifieur $0.964$. Elles tombent du même côté de $\frac12$, si bien que l'erreur ne porte que sur la confiance. C'est l'ajout d'une seconde variable qui en fait une mauvaise décision.

\iqwrong{Dire que le Bayes naïf est insensible à la corrélation parce qu'il classe souvent bien. C'est souvent vrai pour le classement, mais des variables corrélées produisent une surconfiance qui peut modifier les décisions.}

\iqfollow{La régression logistique souffrirait-elle de la même façon~?}{Non. Elle répartit le coefficient entre les copies, de sorte que leur somme reste la même, et elle renvoie le même ajustement. Un essai avec trois copies le confirme.}

\iqwhy{Elle rend concret ce que coûte l'hypothèse d'indépendance. En finance, où les variables sont fortement corrélées, c'est un problème quotidien.}
\end{solution}

\begin{solution}{iq:iv1:machine-learning:21}
Les deux gradients valent $0.5$, les nouveaux poids sont $-0.05$ et $1.95$, et la perte passe de $0.5$ à $0.452$.

Passe avant~: $h = \sigma(0) = 0.5$, $\hat{y} = 2 \times 0.5 = 1$, et la perte vaut $\frac12$.

Rétropropagation\index{rétropropagation}~: $\partial L/\partial\hat{y} = \hat{y} - y = 1$.

Pour $w_2$~: $\partial\hat{y}/\partial w_2 = h = 0.5$, donc le gradient vaut $0.5$.

Pour $w_1$~: $\partial\hat{y}/\partial h = w_2 = 2$, et $\partial h/\partial w_1 = \sigma'(0)x = 0.25$, donc le gradient vaut $1 \times 2 \times 0.25 = 0.5$. Les différences finies centrées concordent avec les deux à huit décimales.

Pas de descente~: $w_1 = 0 - 0.1(0.5) = -0.05$, et $w_2 = 2 - 0.1(0.5) = 1.95$. Puis $h = \sigma(-0.05) = 0.4875$, $\hat{y} = 0.9506$, et la perte vaut $0.4518$. Une estimation au premier ordre de la baisse donne $0.1 \times (0.5^2 + 0.5^2) = 0.05$, et la baisse réelle vaut $0.048$~: les chiffres sont cohérents.

\iqwrong{Oublier la dérivée de la sigmoïde, $\sigma(1 - \sigma)$, qui vaut ici $\frac14$. Ou mettre à jour $w_2$ d'abord, puis utiliser le nouveau $w_2$ dans le gradient de $w_1$. Les deux gradients proviennent de la même passe avant.}

\iqfollow{Pourquoi calculer les dérivées en remontant, et non en avançant~?}{Une seule perte et de nombreux poids~: une passe arrière donne toutes les dérivées partielles pour à peu près le coût d'une passe avant, alors qu'une passe avant demanderait un parcours par poids.}

\iqwhy{Tout candidat qui affirme avoir de l'expérience des réseaux de neurones doit savoir faire cela de tête. La vérification par rapport à la baisse prédite témoigne de la rigueur.}
\end{solution}

\begin{solution}{iq:iv1:machine-learning:22}
La borne vaut $4^{-10} \approx 10^{-6}$. Le ReLU a une dérivée égale à 1 là où il est actif, mais des unités peuvent mourir. La bonne variance est $2/n$.

La dérivée de la sigmoïde vaut au plus $\frac14$, au milieu de son intervalle. Avec des poids unitaires, la règle de la chaîne donne un produit de dix facteurs de ce type, au plus $9.5 \times 10^{-7}$\index{évanescence du gradient}. Les premières couches n'apprennent presque plus. Des poids plus grands ne règlent pas le problème en général, car ils poussent les unités vers la saturation, où la dérivée est proche de zéro.

Le ReLU a une dérivée égale à 1 sur sa partie active, donc le produit ne se réduit pas le long des chemins actifs. Le nouvel échec est qu'une unité négative sur toutes les données a un gradient nul et ne se rétablit jamais~: c'est un neurone mort\index{ReLU morte}. Un taux d'apprentissage élevé en est la cause habituelle.

Pour la variance initiale~: avec des poids de variance $v$ et $n$ entrées de moyenne quadratique $q$, la pré-activation a pour variance $nvq$. Le ReLU conserve la moitié de son second moment. Donc $q_{\text{suivant}} = \frac12 nvq$, qui est constant pour $v = 2/n$\index{initialisation}. Un passage sur 30 couches de largeur 500 le confirme~: le signal reste dans un facteur 10 avec $2/n$, il s'effondre sous $10^{-5}$ avec $1/n$, et il explose au-delà de $10^5$ avec $4/n$.

\iqwrong{Prendre $v = 1/n$ avec le ReLU, ce qui convient à une activation symétrique mais fait perdre la moitié du signal à chaque couche.}

\iqfollow{Que fait une connexion résiduelle pour le gradient~?}{La dérivée de $x + g(x)$ vaut $1 + g'(x)$. Un chemin passant par l'identité transporte donc le gradient sans atténuation, quelle que soit la profondeur du réseau.}

\iqwhy{Elle teste ce que font réellement les remèdes usuels. Une valeur exacte de la variance montre que le candidat sait les dériver et ne se contente pas de les nommer.}
\end{solution}

\begin{solution}{iq:iv1:machine-learning:23}
Aucune droite ne sépare les classes, par une somme de deux inégalités. Le réseau $y = r(x_1 + x_2) - 2r(x_1 + x_2 - 1)$, avec $r$ le ReLU, le calcule.

Supposons $w_1x_1 + w_2x_2 + c > 0$ exactement pour la classe 1. Les entrées $(1, 0)$ et $(0, 1)$ donnent $w_1 + c > 0$ et $w_2 + c > 0$. Les entrées $(0, 0)$ et $(1, 1)$ donnent $c < 0$ et $w_1 + w_2 + c < 0$ (ou des inégalités larges). En additionnant la première paire, on obtient $w_1 + w_2 + 2c > 0$, et en additionnant la seconde, $w_1 + w_2 + 2c < 0$. C'est une contradiction\index{ou exclusif}.

Avec $s = x_1 + x_2$, le réseau donne~:
\begin{itemize}
\item $(0, 0)$~: $s = 0$, donc $r(0) - 2r(-1) = 0$.
\item $(1, 0)$ ou $(0, 1)$~: $s = 1$, donc $r(1) - 2r(0) = 1$.
\item $(1, 1)$~: $s = 2$, donc $r(2) - 2r(1) = 2 - 2 = 0$.
\end{itemize}
Les deux unités cachées créent un coude dans la fonction de $s$, ce qui est précisément ce qui manque à un modèle linéaire. Un arbre de profondeur 2 peut aussi le faire.

\iqwrong{Affirmer qu'une couche cachée rend le problème linéairement séparable sans construction. Ou utiliser deux unités de même seuil, qui restent une fonction linéaire de $s$ au-delà de ce seuil.}

\iqfollow{Comment cela se généralise-t-il à la parité de $n$ bits~?}{Une couche cachée de $n$ unités ReLU suffit~: avec $s$ la somme des bits, $r(s) + \sum_{k = 1}^{n - 1}2(-1)^kr(s - k)$. Une vérification sur toutes les entrées jusqu'à $n = 5$ le confirme.}

\iqwhy{C'est la raison classique des couches cachées. Écrire les poids est ce qu'un entretien ajoute au slogan.}
\end{solution}

\begin{solution}{iq:iv1:machine-learning:24}
L'entropie croisée donne $-0.999$ et la perte quadratique $-0.00091$, soit un facteur d'environ 1100. L'entropie croisée est convexe en $z$, ce qui n'est pas le cas de la perte quadratique.

Ici $p = \sigma(-7) = 0.00091$.

L'entropie croisée $-\ln p$ a pour dérivée $\partial/\partial z = p - 1 = -0.999$\index{entropie croisée}. Ainsi, lorsque le modèle se trompe lourdement, le gradient est proche de sa valeur maximale en valeur absolue.

La perte quadratique a $\partial/\partial z = (p - 1) \times \sigma'(z) = (p - 1)p(1 - p) = -9.09 \times 10^{-4}$. La règle de la chaîne fait intervenir $\sigma'$, qui est minuscule dans les queues. Ainsi, la pire erreur ne produit presque aucun gradient. Le facteur vaut $1/(p(1 - p)) \approx 1100$.

La dérivée seconde de l'entropie croisée en $z$ vaut $p(1 - p) > 0$, elle est donc convexe. La perte quadratique a des différences secondes négatives en $z = -7$ et positives en $z = 1$, elle n'est donc pas convexe en $z$, et elle présente des régions plates d'où l'entraînement progresse à peine.

\iqwrong{Affirmer que la perte quadratique est fausse parce qu'elle n'est pas une vraisemblance. Elle est la vraisemblance d'un autre modèle. Le problème, pour l'entraînement, c'est le gradient.}

\iqfollow{Quels sont les gradients lorsque le modèle a raison, en $z = 7$ et $y = 1$~?}{$-0.00091$ avec l'entropie croisée, et $-8.3 \times 10^{-7}$ avec la perte quadratique. L'entropie croisée cesse de pousser dès que le modèle a raison, et seulement alors.}

\iqwhy{Elle explique un choix de conception que presque tout le monde accepte sans examen, et elle donne un nouvel exemple où une petite dérivée fait des dégâts.}
\end{solution}

\begin{solution}{iq:iv1:machine-learning:25}
L'aire vaut $10/12 = 0.833$ par chacune des deux méthodes.

Trier par score, puis faire descendre le seuil. Chaque positif élève la courbe ROC de $1/n_+$, et chaque négatif la décale vers la droite de $1/n_-$. L'aire sous l'escalier\index{courbe ROC} est une somme sur les négatifs de la hauteur atteinte lorsqu'on passe chacun d'eux, soit $1/n_-$ fois la part des positifs classés au-dessus. L'aire vaut donc
\[
\frac{1}{n_+n_-}\sum_{+, -}\mathbf{1}(s_+ > s_-),
\]
la part des paires bien ordonnées, les ex aequo comptant pour $\frac12$ grâce au segment diagonal.

Paires~: $0.9$ bat les quatre négatifs. $0.6$ bat $0.3$, $0.2$, $0.1$, ce qui fait 3. $0.4$ bat aussi ces trois-là. Le total est de 10 sur 12.

Rangs depuis le bas\index{statistique de Mann-Whitney}~: $0.1, 0.2, 0.3$ prennent 1, 2, 3~; $0.4$ prend 4~; $0.6$ prend 5~; $0.7$ prend 6~; $0.9$ prend 7. La somme des rangs des positifs vaut 16. Alors $U = 16 - 3 \times 4/2 = 10$, le même décompte. Une somme par trapèzes sur la courbe ROC donne aussi $0.833$.

\iqwrong{Affirmer que l'aire est la justesse pour un seuil donné. Elle fait intervenir tous les seuils, et elle porte sur l'ordre des paires, non sur le classement des lignes.}

\iqfollow{Que devient l'aire si l'on prend l'exponentielle de tous les scores~?}{Rien. Elle ne dépend que des rangs, si bien qu'un modèle mal calibré peut avoir une très bonne aire.}

\iqwhy{Le résultat donne un sens à un nombre souvent cité sans en avoir un, et la forme par somme des rangs est celle qu'on utilise en pratique pour le calculer.}
\end{solution}

\begin{solution}{iq:iv1:machine-learning:26}
Il faut agir au-dessus de la probabilité $0.1$ dans le premier cas et au-dessus de $2/3$ dans le second. Avec des scores non calibrés, ces seuils ne veulent pas dire ce qu'ils disent.

Agir lorsque le coût espéré de l'action est inférieur au coût espéré de l'inaction\index{seuil de décision}~: $(1 - p) \times 1 < p \times 9$, ce qui donne $p > 1/10$. En général, $p^* = c_{\text{FP}}/(c_{\text{FP}} + c_{\text{FN}})$. Une erreur coûteuse abaisse le seuil.

Pour la transaction, l'espérance $0.5p - 1(1 - p)$ est positive lorsque $p > 1/1.5 = 2/3$. En $p = 2/3$, l'espérance vaut exactement zéro~: $\frac13 - \frac13$.

Ces règles supposent que le score est une vraie probabilité. Un score de $0.7$ issu d'un modèle trop confiant peut signifier en réalité $0.55$, auquel cas la transaction perd de l'argent tout en paraissant sûre. Il faut donc vérifier la calibration\index{calibration}~: parmi les lignes dont le score est proche de $0.7$, environ 70~\% doivent être positives, sur des données postérieures à la période d'entraînement. On recalibre alors, ou l'on choisit le seuil directement sur une période de validation, d'après le profit mesuré.

\iqwrong{Prendre $\frac12$ par défaut, ce qui n'est juste que si les deux erreurs coûtent la même chose. Ou choisir le seuil qui maximise la justesse alors que les coûts ne sont pas égaux.}

\iqfollow{Que change un coût de $0.1~\%$ par transaction~?}{$0.5p - (1 - p) - 0.1 > 0$ donne $p > 1.1/1.5 = 0.733$. De petits coûts relèvent fortement la barre lorsque le gain est faible.}

\iqwhy{Une prédiction n'est utile qu'à travers une décision. C'est bon signe lorsqu'un candidat traduit les coûts en seuil, et ne s'interroge sur la calibration qu'ensuite.}
\end{solution}

\begin{solution}{iq:iv1:machine-learning:27}
Sharpe quotidien $0.04$, annuel $0.64$. Avec le coût, le chiffre annuel vaut $0.32$. Un coût de $0.04$ le ramène à zéro.

Le rendement quotidien moyen vaut $0.52 - 0.48 = 0.04$\index{taux de réussite}. La variance vaut $1 - 0.04^2$, donc l'écart-type vaut $0.9992$. Le ratio de Sharpe quotidien vaut $0.04$, et le ratio annuel $0.04\sqrt{252} = 0.635$. Une longue simulation le confirme.

Un coût de $0.02$ par transaction divise par deux la moyenne, qui vaut alors $0.02$, et divise par deux le ratio de Sharpe, qui vaut $0.32$. La volatilité ne change pas. Le coût de point mort est donc exactement l'edge, soit $0.04$.

Deux leçons en découlent. Des taux de réussite très modestes comptent~: 52~\% ressemble à une pièce équilibrée et donne un chiffre respectable, si les paris sont indépendants. Et l'edge par transaction est si faible que le coût décide facilement de tout. On ne peut pas juger la précision d'un modèle sans connaître son turnover ni ses coûts.

\iqwrong{Supposer que le ratio de Sharpe issu du seul taux de réussite est robuste. Tout le calcul repose sur des gains et des pertes de même taille, ce qui échoue souvent.}

\iqfollow{Quel taux de réussite permet d'atteindre le point mort si les gains valent $0.9$ fois la taille des pertes~?}{$0.9p = 1 - p$ donne $p = 1/1.9 = 0.526$. Un modèle à 52~\% perd donc de l'argent.}

\iqwhy{Cela transforme les métriques d'apprentissage automatique en métriques de trading, et c'est là qu'un data scientist doit en général convertir les unités pour la première fois.}
\end{solution}

\begin{solution}{iq:iv1:machine-learning:28}
La corrélation vaut $1 - k/20$. On retire 38 lignes, ce qui laisse 762. Un découpage aléatoire affiche une corrélation supérieure à $0.9$ sur du bruit pur.

Des étiquettes séparées de $k$ jours partagent $20 - k$ de leurs 20 rendements quotidiens indépendants, donc leur corrélation vaut $(20 - k)/20$\index{chevauchement des étiquettes}. À 5 jours, elle vaut $0.75$.

Le bloc de test couvre les jours 400 à 599, en numérotation à partir de zéro. L'étiquette du jour $t$ utilise les jours $t + 1$ à $t + 20$. Une étiquette d'entraînement chevauche une étiquette de test si les deux débuts sont séparés de moins de 20 jours. Il faut donc retirer les 19 dernières lignes d'entraînement avant le bloc et les 19 premières après lui. Il reste $800 - 38 = 762$\index{purge}. Compter directement les lignes donne la même réponse.

Un découpage aléatoire à 5 plis place dans l'entraînement des voisins de chaque ligne de test, car une ligne sur cinq seulement est mise de côté. Le plus proche voisin est en général à 1 jour, avec une corrélation de $0.95$. Ainsi, en se fondant sur la seule date, le modèle «~prédit~» du bruit pur avec une corrélation supérieure à $0.9$ dans une simulation. Avec des blocs et une purge, le même modèle donne une corrélation proche de zéro sur de nombreuses exécutions. Rien n'a changé dans le modèle.

\iqwrong{Penser qu'il suffit de retirer les données qui précèdent le bloc de test. Les étiquettes d'entraînement situées après le bloc partagent aussi des jours, puisqu'une étiquette de test va vers l'avant. Ou mélanger les données parce que le modèle n'a pas de variables temporelles. La fuite se trouve dans les étiquettes.}

\iqfollow{Vous laissez aussi un écart de 5 jours après le bloc, au cas où les variables seraient autocorrélées. Combien de lignes restent-elles~?}{757. Cela s'appelle un embargo, et sa longueur vient de la mémoire des variables et non de l'étiquette.}

\iqwhy{C'est l'erreur classique lorsqu'on applique l'apprentissage automatique aux rendements. La réponse attendue contient le décompte, ainsi que le résultat frappant obtenu sur du bruit pur.}
\end{solution}

\begin{solution}{iq:iv1:machine-learning:29}
(a) est correct~; les autres présentent une fuite. (b) comporte un biais d'anticipation dans une variable lente, (c) un biais d'anticipation sur l'échelle, (d) utilise une information qui n'est pas encore publique, (e) est un biais de survie, (f) utilise le futur du jour prédit.

(a) La clôture de la veille est connue avant l'ouverture du jour $t$, c'est donc une variable légitime.

(b) Le secteur actuel d'une entreprise peut refléter ce qu'elle est devenue plus tard~: une société reclassée après le décollage de son activité reporte ce résultat sur tout son passé. L'effet est en général faible, mais réel. Il faut utiliser la classification connue à l'époque.

(c) La moyenne et la variance calculées sur tout l'échantillon contiennent la période de test. Lorsque l'échelle d'une variable dérive\index{biais d'anticipation}, cela informe le modèle sur des valeurs postérieures. Il faut utiliser des estimations sur la période d'entraînement ou des estimations glissantes.

(d) Les résultats d'un trimestre fiscal sont publiés plusieurs semaines après la fin du trimestre. Les placer au dernier jour du trimestre donne au modèle plusieurs semaines d'avance sur l'information. Il faut utiliser la date d'annonce et les premiers chiffres publiés, et non les chiffres retraités ultérieurement.

(e) Les membres actuels de l'indice sont les survivants et les gagnants. L'appartenance renseigne le modèle sur le fait que l'action a suffisamment bien performé pour rester\index{biais de survie}. Il faut utiliser l'appartenance telle qu'elle était à chaque jour.

(f) Le volume du jour $t$ n'est connu qu'à la clôture, et il est lié à l'ampleur du mouvement de la journée. C'est une fuite typique, qui produit un backtest impressionnant.

\iqwrong{Considérer qu'une fuite se réduit à l'usage de l'étiquette. La plupart des fuites tiennent au moment où une valeur était réellement connue, et chaque variable a besoin d'un horodatage.}

\iqfollow{Comment détecter une fuite sans liste comme celle-ci~?}{On se méfie d'un résultat trop beau. Puis on décale chaque variable d'un jour vers l'avant, on observe celles qui perdent presque toute leur valeur, et on entraîne sur des étiquettes mélangées pour vérifier que le pipeline ne produit aucune compétence.}

\iqwhy{Repérer les fuites occupe une grande part du travail réel, et l'exemple au cas discutablement bénin teste autant le jugement que la connaissance.}
\end{solution}

\begin{solution}{iq:iv1:machine-learning:30}
Il reste sur la partition haut/bas, avec une somme des carrés de 16. L'optimum est gauche/droite, avec une somme de 1. Un départ aléatoire tombe sur le mauvais résultat avec une probabilité de $\frac13$, et cinq redémarrages échouent tous avec une probabilité de $0.4~\%$.

Chaque point rejoint le centre le plus proche, si bien que les deux points du haut et les deux du bas forment des groupes. Leurs moyennes sont les mêmes centres $(0, \pm 0.5)$. L'algorithme a convergé\index{k-moyennes}. Chaque point est à 2 de son centre horizontalement, donc la somme des carrés vaut $4 \times 4 = 16$.

La partition gauche/droite a pour centres $(\pm 2, 0)$. Chaque point est à $0.5$ de son centre, donc la somme des carrés vaut $4 \times 0.25 = 1$, ce qui est bien meilleur. Chaque étape de l'algorithme fait diminuer la fonction objectif, mais il n'atteint qu'un minimum local.

Il existe six paires de points de données à partir desquelles démarrer. Prenons la paire du petit côté gauche, $(-2, 0.5)$ et $(-2, -0.5)$. Les points de droite $(2, 0.5)$ et $(2, -0.5)$ se répartissent alors selon leur hauteur, et le mauvais résultat haut/bas apparaît. Il en va de même pour la paire du côté droit. Les quatre autres paires mènent à gauche/droite, comme le confirme le calcul sur les six. La probabilité vaut donc $2/6$, et cinq redémarrages échouent tous avec une probabilité de $(1/3)^5 = 0.004$.

\iqwrong{Affirmer que les $k$-moyennes trouvent l'optimum parce que la fonction objectif décroît à chaque étape. Une fonction objectif décroissante montre qu'il converge, et non vers quoi il converge.}

\iqfollow{Un démarrage plus astucieux tire le premier centre au hasard, puis le suivant avec une probabilité proportionnelle au carré de la distance. Depuis $(2, 0.5)$, quelle est la probabilité de la mauvaise paire~?}{$1/(1 + 16 + 17) = 1/34$, soit environ 3~\%. Le carré de la distance favorise les points éloignés, et cette règle est le démarrage $k$-means++.}

\iqwhy{Un algorithme qui s'arrête n'est pas un algorithme qui réussit. Cet exemple réduit donne une explication exacte de l'usage des redémarrages et des démarrages soignés.}
\end{solution}

\begin{solution}{iq:iv1:machine-learning:31}
La validation croisée affiche une erreur inférieure à 5~\%, proche de zéro, alors que l'erreur réelle vaut 50~\%. La sélection de variables doit se faire à l'intérieur de chaque pli.

Avec 5000 variables de bruit et seulement 50 lignes, certaines variables sont corrélées aux étiquettes par hasard. Les 100 meilleures sont sélectionnées à partir de toutes les lignes, y compris celles qui seront mises de côté plus tard. Ainsi, dans chaque pli, les variables connaissent déjà les étiquettes mises de côté. Un classifieur au centroïde le plus proche a alors une tâche facile\index{fuite liée à la sélection de variables}.

La bonne procédure traite la sélection comme une partie du modèle. Dans chaque pli, on sélectionne les 100 variables en n'utilisant que les lignes d'entraînement, on ajuste sur ces lignes, puis on prédit les lignes mises de côté. Désormais, les étiquettes des lignes mises de côté n'ont aucun effet sur les variables, et l'erreur vaut environ 50~\%, comme il se doit, puisque les étiquettes sont aléatoires.

Une simulation portant sur 40 jeux de données aléatoires donne une moyenne inférieure à 5~\% avec la mauvaise procédure, et $50~\%$, à bruit près, avec la bonne. Rien d'autre ne change.

La même erreur apparaît dès qu'un prétraitement dépendant de la cible, de toute nature, est effectué en premier~: codage de la cible (target encoding), rééchantillonnage, choix de l'univers d'après les performances ultérieures, ou choix de la plage de réglage en regardant les résultats sur tout l'échantillon.

\iqwrong{Penser que, puisque le classifieur n'est entraîné que sur les plis d'entraînement, la validation est honnête. Le classifieur n'est que la dernière étape. Tout ce qui a vu les étiquettes compte.}

\iqfollow{Est-ce aussi une fuite que de retirer les variables à faible variance avant la validation croisée~?}{Non, du point de vue des étiquettes, puisque ce filtre ne les utilise pas. À terme, en revanche, utiliser la variance de tout l'échantillon divulgue une part d'information sur l'échelle.}

\iqwhy{C'est une erreur très répandue dans le domaine, au résultat spectaculaire. On attend du candidat qu'il dise d'emblée où doit se placer le filtre.}
\end{solution}

\begin{solution}{iq:iv1:machine-learning:32}
On mesure environ $2.80$, et le ratio de Sharpe réel espéré vaut $0.56$. Avec quatre fois l'historique, les chiffres sont $1.77$ et $0.89$.

Les valeurs mesurées suivent une loi normale de variance $\tau^2 + \sigma^2 = 0.25 + 1 = 1.25$. Le plus grand de 100 normales centrées réduites indépendantes vaut $2.51$ en moyenne, donc la meilleure valeur mesurée vaut $2.51 \times \sqrt{1.25} = 2.80$.

Étant donnée une valeur mesurée $m$, la valeur réelle a pour espérance $\frac{\tau^2}{\tau^2 + \sigma^2}m = 0.2m$\index{biais de sélection}. La sélection ne dépend que des valeurs mesurées, donc cela s'applique aussi au gagnant, et $0.2 \times 2.80 = 0.56$. Ainsi, 80~\% du ratio de Sharpe mesuré est du bruit, et ce bruit ne reviendra pas. Une simulation de 200\,000 de telles compétitions donne $2.80$ et $0.56$.

Avec quatre fois l'historique, l'écart-type du bruit se divise par deux, pour valoir $0.5$. La variance de la mesure vaut désormais $0.5$, et la meilleure valeur mesurée vaut $2.51 \times 0.707 = 1.77$. Le facteur de rétrécissement vaut $\frac12$, ce qui donne $0.89$. Le chiffre mesuré baisse donc, et le chiffre réel monte~: de meilleures preuves donnent un nombre plus petit et plus réaliste.

À la limite d'une mesure très bruitée, le facteur tend vers zéro, et le gagnant ne vaut pas mieux que le candidat moyen, comme il se doit. Avec une mesure parfaite, le facteur vaut 1, et la valeur réelle du gagnant vaut $0.5 \times 2.51 = 1.25$.

\iqwrong{Diviser par $\sqrt{100}$, ou par le nombre d'essais, en guise de règle empirique. La bonne décote dépend de l'a priori et de la quantité de données, et le nombre d'essais n'intervient que par la taille du maximum.}

\iqfollow{Si les ratios de Sharpe réels avaient une moyenne a priori de $-0.2$, après coûts, quelle est la valeur réelle espérée du gagnant~?}{Environ $-0.2 + 0.2 \times 2.80 = 0.36$. Une moyenne a priori négative se retranche en entier, car les écarts mesurés portent sur la moyenne.}

\iqwhy{Un modèle compact de la recherche~: beaucoup d'idées, des tests bruités et une malédiction du vainqueur. Il fournit un chiffre défendable pour le degré de confiance à accorder à un backtest.}
\end{solution}

\begin{solution}{iq:iv1:machine-learning:33}
L'erreur vaut $1/w + q(w + 1)(2w + 1)/(6w)$, avec un minimum en $w \approx \sqrt{3/q} = 173$, où elle vaut $0.0116$. La courbe est plate près du minimum.

La prévision utilise la moyenne des $w$ dernières observations pour prévoir le niveau un pas plus loin. Le bruit d'observation contribue pour $\sigma^2/w = 1/w$. Pour le niveau, l'erreur de prévision est une somme pondérée de pas d'une marche aléatoire\index{longueur de fenêtre}. Le pas le plus récent est entièrement manqué, avec un poids de 1, le précédent avec un poids de $(w - 1)/w$, et ainsi de suite jusqu'à $1/w$. La variance vaut
\[
q\sum_{k = 1}^{w}(k/w)^2 = q\frac{(w + 1)(2w + 1)}{6w}.
\]

Pour $w$ grand, la somme des deux termes vaut environ $1/w + qw/3$, avec un minimum en $w = \sqrt{3/q} = 173$. Une recherche sur les entiers donne 173, et une simulation donne la même erreur, $0.0116$.

La sensibilité est faible. Pour $w = 100$ ou $w = 300$, l'erreur n'augmente que de 15~\%. Une fenêtre peut donc être décalée d'un facteur proche de 2 sans grand dommage, et les deux cas limites sont clairs~: lorsque $q \to 0$, il faut utiliser toutes les données, et lorsque $q$ augmente, la fenêtre rétrécit.

La leçon pour les fenêtres d'entraînement est que la bonne longueur est fixée par la vitesse à laquelle la relation dérive par rapport à son bruit. Comme aucun des deux n'est bien connu, régler ce paramètre finement par backtest risque fort de conduire à du surapprentissage.

\iqwrong{Affirmer que davantage de données est toujours préférable, ou choisir la fenêtre qui a donné les meilleurs résultats historiques à plusieurs décimales près. Le minimum plat signifie que le choix ne peut pas être déterminé avec précision, et qu'il n'a pas besoin de l'être.}

\iqfollow{Qu'apporte ici la pondération exponentielle~?}{Le filtre optimal a un gain d'environ $\sqrt{q} = 0.01$, ce qui correspond à une demi-vie de 69 pas, et une erreur de $0.01005$, inférieure de 13~\% à celle de la meilleure fenêtre. Il atténue progressivement le poids des anciennes données, et c'est le filtre de Kalman en régime permanent pour ce modèle.}

\iqwhy{Tout modèle appliqué à des données non stationnaires doit choisir une fenêtre. Peu de candidats savent en établir une et dire quelle importance a ce choix.}
\end{solution}

\begin{solution}{iq:iv1:machine-learning:34}
L'aire vaut $\Phi(d/\sqrt{2})$, ce qui donne $0.528$, $0.638$ et $0.760$. Une aire de $0.55$ correspond à $d = 0.178$, et à une corrélation de $0.089$ avec l'étiquette.

L'aire est la probabilité qu'un positif reçoive un score supérieur à celui d'un négatif\index{courbe ROC}. La différence entre le score du positif et celui du négatif suit une loi normale, de moyenne $d$ et de variance 2. La probabilité qu'elle soit positive vaut donc $\Phi(d/\sqrt{2})$.

Pour la seconde partie, on inverse~: $d = \sqrt{2}\,\Phi^{-1}(0.55) = \sqrt{2} \times 0.1257 = 0.178$.

Pour des classes équilibrées, l'étiquette $y \in \{0, 1\}$ a une variance de $\frac14$. Le score a pour variance $1 + d^2/4$, car le mélange des moyennes ajoute $d^2/4$. La covariance entre les deux vaut $d/4$. Donc
\[
\operatorname{corr} = \frac{d/4}{\frac12\sqrt{1 + d^2/4}} = \frac{d/2}{\sqrt{1 + d^2/4}} = 0.089.
\]
Une simulation confirme à la fois $0.089$ et l'aire de $0.55$. Pour $d = 0$, l'aire vaut $\frac12$ et la corrélation est nulle, comme il se doit.

Pour la prévision des rendements, une corrélation de $0.09$ avec le signe du rendement est un effet important. Ainsi, une aire de $0.55$, qui paraît faible à qui est habitué aux images et au texte, est très forte si elle est réelle, et elle mérite la méfiance.

\iqwrong{Utiliser les repères d'autres domaines, où $0.55$ est un échec. Ou assimiler la corrélation à l'écart de l'aire par rapport à $\frac12$.}

\iqfollow{Quelle précision ce modèle atteint-il au seuil situé au milieu~?}{$\Phi(d/2) = 0.535$. Une précision directionnelle de $53.5~\%$ correspond à l'aire de $0.55$.}

\iqwhy{Convertir entre métriques est la façon dont les résultats se comparent d'une équipe à l'autre. Cela sert aussi à situer ce qui paraît impressionnant en finance.}
\end{solution}

\begin{solution}{iq:iv1:machine-learning:35}
Les probabilités corrigées valent $0.091$ et $0.474$. Dans une régression logistique, la correction ne modifie que l'ordonnée à l'origine, de $\ln 0.1 = -2.30$, et non les pentes.

Soit $\beta = 0.1$ la part des négatifs conservés. La conservation est indépendante des variables, à classe donnée. Donc la formule de Bayes donne\index{sous-échantillonnage}
\[
\frac{p'}{1 - p'} = \frac{P(y = 1 \mid x)}{\beta\,P(y = 0 \mid x)} = \frac{1}{\beta}\cdot\frac{p}{1 - p}.
\]
Ainsi, la cote sur les données échantillonnées vaut dix fois la cote réelle, quel que soit $x$, et la correction s'écrit $p = \beta p'/(\beta p' + 1 - p')$.

Pour $p' = 0.5$, la cote échantillonnée vaut 1, et la cote réelle vaut $0.1$, donc $p = 0.1/1.1 = 0.091$. Pour $p' = 0.9$, la cote échantillonnée vaut 9, et la cote réelle vaut $0.9$, donc $p = 0.9/1.9 = 0.474$. Une observation que le modèle évalue à $90~\%$ relève du pile ou face.

Dans une régression logistique, les log-cotes sont linéaires en $x$, et l'échantillonnage les décale d'une constante, $-\ln\beta$. Ainsi l'ajustement échantillonné a les mêmes pentes, avec une ordonnée à l'origine supérieure de $\ln 10 = 2.30$, et la correction retranche cette quantité. Une simulation avec 400\,000 lignes, une ordonnée à l'origine de $-3$ et une pente de 1 donne une pente de $1.0$ et une ordonnée à l'origine de $-0.70$, comme prévu.

Ce résultat repose sur le fait que la conservation est indépendante de $x$, à classe donnée. Un filtre qui dépend aussi des variables la brise, et le décalage n'est alors plus constant. Les métriques de classement ne changent pas sous cette correction, alors que les probabilités changent beaucoup. Toute règle de décision fondée sur elles doit donc utiliser les valeurs corrigées.

\iqwrong{Utiliser directement les scores des données équilibrées comme probabilités dans une règle de coût. Ou appliquer une même mise à l'échelle à la probabilité, alors qu'il faut corriger la cote.}

\iqfollow{Quelle est une alternative à la correction a posteriori~?}{On pondère chaque négatif conservé par 10 dans la perte. L'ajustement pondéré estime directement la population d'origine, sans correction, mais les estimations sont plus bruitées qu'avec toutes les données.}

\iqwhy{Le sous-échantillonnage est courant pour les événements rares. Beaucoup de chaînes de traitement produisent des nombres mal calibrés, parce que cette étape est oubliée.}
\end{solution}

\begin{solution}{iq:iv1:machine-learning:36}
Les log-cotes initiales valent $\ln 1.5 = 0.405$. La valeur de feuille issue de Newton vaut $0.625$, l'optimum exact vaut $0.693$, et la nouvelle probabilité vaut $0.615$.

Avec $F$ les log-cotes et $p = \sigma(F)$, la perte sur une observation vaut $-y\ln p - (1 - y)\ln(1 - p)$. Son gradient en $F$ est $p - y$, et sa hessienne est $p(1 - p)$\index{perte logarithmique}. La meilleure constante est le taux de base, donc $F_0 = \ln(0.6/0.4)$.

Pour $p = 0.6$, un positif a pour gradient $-0.4$, un négatif a $+0.6$, et chaque observation a pour hessienne $0.24$.

Pour la feuille~: $G = 30(-0.4) + 10(0.6) = -6$, et $H = 40(0.24) = 9.6$. Le pas de Newton vaut $-G/H = 0.625$.

L'optimum exact rend la probabilité de la feuille égale à sa part de positifs, $\frac34$, dont les log-cotes valent $\ln 3 = 1.099$. Ainsi le pas exact vaut $1.099 - 0.405 = 0.693$, et une recherche directe le confirme. Ici, Newton reste environ 10~\% en deçà, et il s'améliore aux itérations suivantes.

Avec un rétrécissement de $0.1$, la mise à jour vaut $F = 0.405 + 0.0625 = 0.468$, donc $p = 0.615$, en hausse par rapport à $0.6$, en direction de la valeur $0.75$ de la feuille.

\iqwrong{Prendre la moyenne brute des résidus, $6/40 = 0.15$, comme valeur de feuille. C'est un pas sur l'échelle des probabilités, et non sur celle des log-cotes. Ou oublier la hessienne, qui sert justement à passer de l'une à l'autre.}

\iqfollow{L'autre feuille compte 30 positifs et 30 négatifs. Quelles sont sa valeur de Newton et sa valeur exacte~?}{Newton donne $-6/14.4 = -0.417$. Le pas exact vaut $\ln 1 - 0.405 = -0.405$. Ici, Newton dépasse légèrement l'optimum.}

\iqwhy{Le boosting du second ordre est l'outil standard sur données tabulaires, et peu de candidats savent calculer un pas à la main. Le faire permet de savoir ce qui est réellement ajusté dans une feuille.}
\end{solution}

\begin{solution}{iq:iv1:machine-learning:37}
$w^* = -G/(H + \lambda)$, avec gain $\frac12\bigl(\frac{G_L^2}{H_L + \lambda} + \frac{G_R^2}{H_R + \lambda} - \frac{(G_L + G_R)^2}{H_L + H_R + \lambda}\bigr)$. Les gains valent $4.76$ et $26.25$, donc la bibliothèque préfère la coupure déséquilibrée. Un seuil minimal sur la hessienne dans chaque nœud fils protège contre ce cas.

Minimiser la forme quadratique en $w$ donne $w^* = -G/(H + \lambda)$, et le minimum vaut $-\frac12G^2/(H + \lambda)$\index{gain de coupure}. Le gain d'une coupure est le minimum du nœud parent moins ceux des fils, ce qui correspond à la formule ci-dessus. On peut aussi retrancher une pénalité constante par feuille.

Dans les deux cas, le nœud parent a $G = 0$, donc le dernier terme s'annule.

Première coupure~: $\frac12(100/21 + 100/21) = 4.76$, avec des valeurs de feuille $\pm 0.476$.

Deuxième coupure~: $\frac12(100/2 + 100/40) = 26.25$, avec une valeur de feuille gauche de $10/2 = 5$.

La deuxième coupure l'emporte donc largement, alors qu'elle repose sur une hessienne de 1, ce qui apporte très peu d'information~: quelques lignes au plus. Sa valeur de feuille est plus de dix fois supérieure. Une recherche sur la forme quadratique confirme la valeur 5 de la feuille.

Un $\lambda$ plus grand aide, mais même $\lambda = 10$ laisse subsister la préférence~: les gains valent $3.33$ et $5.57$. La protection directe est un seuil minimal sur la somme des hessiennes par nœud fils\index{poids minimal par nœud fils}. Un seuil de 5 rejette tout net la deuxième coupure. Pour $\lambda = 0$, le second gain vaut $51.28$, donc le terme de régularisation fait déjà beaucoup de travail.

\iqwrong{Penser que $\lambda$ n'est qu'une pénalité sur la taille des feuilles, sans effet sur le choix des coupures. Il figure pourtant dans le gain, comme on l'a montré.}

\iqfollow{Pourquoi la hessienne mesure-t-elle l'information~?}{Avec la perte quadratique, chaque ligne a une hessienne de 1, donc $H$ est le nombre de lignes. Avec la perte logarithmique, une ligne a pour hessienne $p(1 - p)$, si bien que les lignes dont le modèle est déjà sûr ne comptent presque pas.}

\iqwhy{C'est la formule utilisée dans les bibliothèques standard, et un paramètre que l'on règle souvent par habitude. Les candidats capables de dériver le gain savent ce qu'ils règlent.}
\end{solution}

\begin{solution}{iq:iv1:machine-learning:38}
Le résultat découle de la formule de Sherman--Morrison, et le résidu leave-one-out du dernier point vaut $-10$ par les deux méthodes.

Soit $A = X^\top X$ et soit $x_i$ la ligne $i$. Retirer la ligne donne $A_{-i} = A - x_ix_i^\top$. La formule de Sherman--Morrison\index{formule de Sherman-Morrison} donne
\[
A_{-i}^{-1} = A^{-1} + \frac{A^{-1}x_ix_i^\top A^{-1}}{1 - h_{ii}}, \qquad h_{ii} = x_i^\top A^{-1}x_i.
\]
Puis $\hat{\beta}_{-i} = A_{-i}^{-1}(X^\top y - x_iy_i)$. En développant et en simplifiant,
\[
\hat{\beta}_{-i} = \hat{\beta} - \frac{A^{-1}x_ie_i}{1 - h_{ii}}.
\]
Ainsi $y_i - x_i^\top\hat{\beta}_{-i} = e_i + h_{ii}e_i/(1 - h_{ii}) = e_i/(1 - h_{ii})$\index{leave-one-out}. Un seul ajustement suffit.

Sur les données~: $\bar{x} = 3.2$, $\sum(x - \bar{x})^2 = 62.8$, et la pente vaut $-5.2/62.8 = -0.083$. La valeur ajustée en 10 vaut $0.637$, donc $e_5 = -0.637$. Le levier vaut $\frac15 + 46.24/62.8 = 0.9363$. La formule donne $-0.637/0.0637 = -10$.

Par son sens~: sans le dernier point, les quatre autres se trouvent exactement sur $y = x$. Cette droite prédit 10 en $x = 10$, alors que $y = 0$. Le résidu vaut donc $0 - 10 = -10$, ce qui concorde. Les ajustements par force brute coïncident sur les cinq lignes.

Le résidu ordinaire paraît petit parce que ce point attire la droite vers lui-même. Son levier vaut $0.94$, sur une somme de $2$ pour les cinq points.

\iqwrong{Juger les valeurs aberrantes à partir des résidus ordinaires. Un point à fort levier peut avoir un résidu minuscule, et c'est pourquoi le rapport compte.}

\iqfollow{La validation croisée généralisée remplace $h_{ii}$ par sa moyenne $p/n$. Que donne-t-elle pour la dernière ligne~?}{$-0.637/(1 - 2/5) = -1.06$, soit un dixième de la taille réelle. C'est une bonne estimation de la moyenne, mais elle masque précisément ce type de point.}

\iqwhy{C'est un résultat élégant, d'usage immédiat, puisque la validation croisée pour les modèles linéaires devient alors gratuite. Le petit exemple montre à quel point un bon résidu peut masquer des choses.}
\end{solution}

\begin{solution}{iq:iv1:machine-learning:39}
La perte espérée est la perte quadratique ordinaire plus $\frac{1 - q}{q}\sum_j\beta_j^2\sum_ix_{ij}^2$. Dans l'exemple, les coefficients valent $q$ fois ceux des moindres carrés~: la moitié pour $q = 0.5$, et $0.8$ d'entre eux pour $q = 0.8$.

Pour une ligne, posons $\tilde{y}_i = \sum_j x_{ij}\beta_j\delta_{ij}/q$. Comme $\mathbb{E}[\delta_{ij}/q] = 1$, la moyenne de $\tilde{y}_i$ vaut $x_i^\top\beta$, la prédiction ordinaire. Donc
\[
\mathbb{E}(y_i - \tilde{y}_i)^2 = (y_i - x_i^\top\beta)^2 + \operatorname{Var}(\tilde{y}_i),
\]
Il s'agit de la décomposition biais--variance pour une seule prédiction. La variance de $\delta/q$ vaut $(1 - q)/q$, et les $\delta_{ij}$ sont indépendants, donc la variance vaut $\frac{1 - q}{q}\sum_jx_{ij}^2\beta_j^2$. En sommant sur les lignes, on obtient la pénalité annoncée\index{dropout}. Une vérification par échantillonnage de la perte espérée, avec 200\,000 masques, concorde à 1~\% près.

Pour des variables non corrélées avec $\sum_ix_{ij}^2 = n$, la pénalité vaut $\lambda\|\beta\|^2$ avec $\lambda = n(1 - q)/q$. La solution est donc
\[
\hat{\beta}_j = \frac{x_j^\top y}{n + \lambda} = \frac{n}{n + n(1 - q)/q}\hat{\beta}^{\text{OLS}}_j = q\,\hat{\beta}^{\text{OLS}}_j.
\]
Pour $q = 1$, rien ne change, comme attendu, et lorsque $q \to 0$, les coefficients tendent vers zéro.

Notons que la pénalité sur chaque coefficient est pondérée par la taille de sa colonne. Le dropout n'est donc pas une régression ridge avec un unique $\lambda$, sauf si les colonnes sont standardisées, et cette pondération le rend insensible aux unités d'une colonne.

\iqwrong{Affirmer que le dropout réduit la sortie du modèle d'un facteur $q$ au moment du test. Avec ce choix d'échelle, la moyenne reste inchangée, et la contraction provient de la variance, via la pénalité.}

\iqfollow{Que dit cela du dropout sur les variables rares~?}{Une variable presque toujours nulle a une somme $\sum x_{ij}^2$ faible, donc son coefficient est peu pénalisé. Contrairement à la régression ridge classique, le dropout ménage les variables rares mais fortes lorsqu'elles sont présentes.}

\iqwhy{Cela transforme une technique qui ressemble à un bricolage en un calcul. L'étape biais--variance est courte, et c'est ce qu'un bon candidat utilise.}
\end{solution}

\begin{solution}{iq:iv1:machine-learning:40}
Un réseau peu profond à $n$ unités a au plus $n + 1$ morceaux. L'application en tente itérée $d$ fois possède $2^d$ morceaux et utilise $2d$ unités. Pour $d = 10$, cela représente 20 unités en profondeur, contre 1023 unités pour le réseau peu profond.

Chaque unité ReLU $r(a x + b)$ a un seul coude, et la sortie est une somme de ces unités, plus une constante. Il y a donc au plus $n$ coudes, et $n + 1$ morceaux\index{profondeur et expressivité}.

L'application en tente sur $[0, 1]$ s'écrit $T(x) = 2r(x) - 4r(x - \frac12)$. Vérification~: en dessous de $\frac12$, le second terme s'annule, et l'on obtient $2x$. Au-dessus de $\frac12$, on obtient $2x - 4x + 2 = 2 - 2x$. Il faut donc deux unités par couche.

Par composition, $T$ envoie chacune de ses deux parties sur tout $[0, 1]$, donc une fonction à $m$ morceaux linéaires, qui couvrent chacun $[0, 1]$, devient une fonction à $2m$ morceaux. Ainsi $T^{\circ d}$ possède $2^d$ morceaux, soit un signal en dents de scie à $2^{d - 1}$ dents. Un décompte des changements de pente sur une grille fine le confirme jusqu'à $d = 10$.

Pour reproduire cela exactement, un réseau peu profond a besoin d'au moins $2^d - 1$ coudes. Pour $d = 10$, cela fait 1023 unités, contre 20. La profondeur procure une économie exponentielle, sur cette fonction. Elle ne dit pas que les fonctions typiques requièrent la profondeur, ni que l'entraînement trouvera de tels poids.

\iqwrong{Affirmer que l'approximation universelle avec une seule couche rend la profondeur inutile. L'approximation universelle ne dit rien sur le nombre d'unités nécessaires.}

\iqfollow{Et pour $d = 30$~?}{Soixante unités en profondeur, contre $2^{30} - 1$, soit plus d'un milliard d'unités en une seule couche.}

\iqwhy{C'est la démonstration la plus nette de l'intérêt de la profondeur, et elle peut se faire à la main en quelques minutes. Elle teste aussi si le candidat connaît les limites de l'affirmation.}
\end{solution}

\begin{solution}{iq:iv1:machine-learning:41}
Le seuil vaut environ $0.11$, proche de la règle empirique $\sqrt{2\ln p/n} = 0.118$. Sur du pur bruit, tout $\lambda$ inférieur laisse entrer des variables de bruit, et avec des signaux faibles, il peut masquer les vraies.

En $\beta = 0$, la partie lisse de la fonction objectif a pour gradient $-X^\top y/n$. Le sous-gradient de $\lambda|\beta_j|$ en 0 peut prendre n'importe quelle valeur dans $[-\lambda, \lambda]$. Donc $\beta = 0$ est optimal, par convexité, exactement lorsque $|x_j^\top y|/n \le \lambda$ pour tout $j$\index{lasso}. On obtient ainsi le seuil, noté $\lambda_{\max}$.

Sur du pur bruit, chaque $x_j^\top y/n$ suit approximativement une loi normale d'écart type $1/\sqrt{n}$. La plus grande valeur absolue parmi 1000 normales de ce type vaut en moyenne $3.43$ écarts types. Donc $\lambda_{\max} \approx 3.43/\sqrt{1000} = 0.108$. L'enveloppe supérieure $\sqrt{2\ln p/n} = 0.118$ surestime légèrement cette valeur pour $p = 1000$. Une simulation donne une valeur dans cet intervalle, un ajustement entièrement nul juste au-dessus, quelques coefficients non nuls à $0.9\lambda_{\max}$, et plus de 50 coefficients non nuls à la moitié de ce seuil, sur du pur bruit.

Ce niveau constitue donc le seuil au-dessus duquel le bruit cesse d'entrer\index{seuil universel}. Un vrai signal n'est détecté que si $|x_j^\top y|/n$ dépasse ce seuil. Avec un coefficient de $0.05$ dans ce cadre, une variable réelle contribue environ $0.05$, ce qui reste sous le maximum du bruit. Donc, à ce niveau, le lasso ne trouve rien, et un $\lambda$ plus faible retient du bruit en même temps que du signal. Avec des signaux faibles répartis sur de nombreuses variables, il ne peut pas les séparer, et la régression ridge prédit souvent mieux.

\iqwrong{Lire les variables sélectionnées comme les vraies. Sur du pur bruit, la validation croisée peut choisir un $\lambda$ qui conserve plusieurs variables.}

\iqfollow{Comment le seuil évolue-t-il si $p$ atteint $10^6$~?}{Seulement en $\sqrt{\ln p}$. L'enveloppe passe de $3.72$ à $5.26$ écarts types, soit une hausse d'environ 40~\%. C'est pourquoi cribler de nombreuses variables coûte moins cher qu'il n'y paraît, du moins au regard de ce seul nombre.}

\iqwhy{La démonstration est courte et donne un nombre d'usage immédiat~: quelle doit être la taille d'un signal pour qu'une méthode parcimonieuse le voie.}
\end{solution}

\begin{solution}{iq:iv1:machine-learning:42}
L'importance par impureté répartit le mérite de la variable réelle entre les deux copies et accorde une part importante à l'identifiant de ligne, pourtant inutile. L'importance par permutation, calculée sur données hors échantillon, donne à l'identifiant une valeur proche de zéro, et sous-estime chaque copie. Il faut permuter ensemble les variables corrélées.

L'importance par impureté additionne les baisses d'impureté\index{importance des variables} dues aux découpages sur chaque variable, mesurées sur les lignes d'entraînement. Avec deux copies, un arbre choisit l'une ou l'autre à chaque découpage, donc le mérite se répartit à peu près également. Chaque copie paraît valoir environ la moitié.

Une variable qui prend une valeur différente sur chaque ligne offre de nombreux découpages possibles. Les arbres profonds peuvent s'en servir pour isoler les lignes bruitées de l'entraînement, et chacun de ces découpages abaisse l'impureté d'entraînement. Elle obtient donc une importance réelle, même sans aucune information. Dans une simulation, elle prend plus de 10~\% du total.

L'importance par permutation mesure la chute du score sur des données hors échantillon lorsqu'une colonne est mélangée\index{importance par permutation}. L'identifiant de ligne n'a aucune valeur sur de nouvelles données, donc la chute est proche de zéro. C'est la bonne lecture.

Mais lorsqu'une copie est mélangée, la forêt peut toujours utiliser l'autre, si bien qu'une variable réelle peut paraître moins importante qu'elle ne l'est. Une simulation montre que la valeur par copie tombe sous les trois quarts de la valeur de la variable seule.

Le remède, pour les variables corrélées, consiste à permuter un groupe ensemble. Dans la simulation, cela restitue plus de 80~\% de la valeur de la variable seule. Retirer le groupe et réajuster le modèle donne une réponse plus coûteuse, mais plus nette. Pour les usages financiers, il faut le faire sur des données ultérieures, et présenter le résultat par groupes, et non variable par variable.

\iqwrong{Écarter une variable parce que son importance est faible, alors qu'elle a un proche parent. Un signal fort partagé entre plusieurs colonnes peut ressembler à plusieurs signaux faibles.}

\iqfollow{Que ne dit pas l'importance~?}{Elle ne dit pas que la variable cause quoi que ce soit, ni qu'elle durera, ni le sens de son effet. Elle décrit la façon dont ce modèle utilise cette colonne sur cet échantillon.}

\iqwhy{La plupart des candidats impriment ces graphiques. Peu savent dire quand ils trompent, ni ce qu'il faut faire à ce sujet.}
\end{solution}

\begin{solution}{iq:iv1:machine-learning:43}
L'erreur de test vaut $1 + p/(n - p - 1)$ pour $p < n - 1$, et $1 + n/(p - n - 1)$ pour $p > n + 1$. Les valeurs sont $2.02$, $11$, $12.1$, $2.01$ et $1.11$.

La nouvelle ligne $x$ est indépendante de l'ajustement, et $y$ a une variance de 1 et n'est pas lié à $x$. L'erreur de test vaut alors $1 + \mathbb{E}\|\hat{\beta}\|^2$, puisque $\mathbb{E}(x^\top\hat{\beta})^2 = \|\hat{\beta}\|^2$ pour un $x$ gaussien standard.

Pour $p < n$~: $\hat{\beta} = (X^\top X)^{-1}X^\top e$, donc, en utilisant $e \perp X$, $\mathbb{E}\|\hat{\beta}\|^2 = \operatorname{tr}\mathbb{E}(X^\top X)^{-1}$. $X^\top X$ suit une loi de Wishart, et la moyenne de son inverse vaut $I/(n - p - 1)$\index{Wishart inverse}. Le terme vaut donc $p/(n - p - 1)$.

Pour $p > n$~: la solution de norme minimale est $\hat{\beta} = X^\top(XX^\top)^{-1}e$, donc $\|\hat{\beta}\|^2 = e^\top(XX^\top)^{-1}e$. Ici $XX^\top$ suit une loi de Wishart de taille $n$ à $p$ degrés de liberté, donc le terme vaut $n/(p - n - 1)$\index{double descente}.

En $p = 50$~: $1 + 50/49$. En $90$~: $1 + 90/9 = 11$. En $110$~: $1 + 100/9 = 12.1$. En $200$~: $1 + 100/99 = 2.01$. En $1000$~: $1 + 100/899 = 1.11$. Les simulations concordent à quelques pour cent près, et à 10~\% près en $p = 90$, où les queues sont épaisses.

Le pic se situe en $p \approx n$. Il existe alors exactement une solution passant par les $n$ points, et l'ajustement de pur bruit exige des coefficients énormes, car la matrice est presque singulière. Au-delà de $n$, il existe de nombreux interpolateurs, et en prendre le plus petit répartit le bruit sur de nombreuses directions, chacune avec un petit coefficient. Lorsque $p \to \infty$, $\hat{\beta} \to 0$ en norme, et l'erreur tend vers 1, le niveau du bruit.

\iqwrong{Conclure que davantage de variables sont gratuites. Ce calcul ne comporte aucun signal. Si le vrai $\beta$ a une taille $r$, répartie uniformément, un terme $r^2(1 - n/p)$ apparaît pour $p > n$, puisque l'ajustement ne voit que la partie de dimension $n$.}

\iqfollow{Avec une taille de signal $r = 1$ et sans bruit, quel est ce terme en $p = 200$~?}{$1 \times (1 - 100/200) = 0.5$. Une simulation donne $0.50$. La moitié de l'énergie du signal est donc simplement hors de portée.}

\iqwhy{C'est la forme exacte la plus simple d'un résultat qui a bouleversé l'image habituelle. Le candidat doit le dériver et éviter les deux conclusions extrêmes.}
\end{solution}

\begin{solution}{iq:iv1:machine-learning:44}
L'erreur vaut environ $\Phi\bigl(-2/\sqrt{4 + 2p/n}\bigr)$, ce qui donne $0.160$, $0.240$ et $0.382$. L'erreur avec moyennes connues vaut $\Phi(-1) = 0.159$.

Avec moyennes connues, la règle est optimale, et dans la direction qui relie les deux moyennes, chaque classe suit une loi normale de variance 1, à la distance $\Delta/2 = 1$ du milieu. L'erreur vaut donc $\Phi(-\Delta/2)$\index{discriminant linéaire}.

Avec moyennes estimées, posons $w = \hat{\mu}_1 - \hat{\mu}_0 = \delta + \varepsilon$, où $\delta$ est la vraie différence et $\varepsilon$ a des composantes indépendantes de variance $2/n$. On classe un nouveau $x$ issu de la classe 1 selon le signe de $w^\top(x - \hat{m})$, où $\hat{m}$ est le milieu des estimations.

Conditionnellement aux estimations, le score suit une loi normale de variance $\|w\|^2$ et d'espérance $w^\top(\mu_1 - \hat{m})$. Or $\mu_1 - \hat{m} = \delta/2 - \eta$, avec $\eta$ un petit bruit, indépendant de $\varepsilon$. Au premier ordre, l'espérance vaut $w^\top\delta/2 \approx \Delta^2/2$, car les termes croisés s'annulent en moyenne et sont petits. Et $\|w\|^2 \approx \Delta^2 + 2p/n$, puisque chacune des $p$ composantes de bruit ajoute $2/n$.

L'erreur vaut donc
\[
\Phi\left(-\frac{\Delta^2/2}{\sqrt{\Delta^2 + 2p/n}}\right) = \Phi\left(-\frac{2}{\sqrt{4 + 2p/n}}\right),
\]
avec $\Delta = 2$ et $n = 50$. Les valeurs sont celles annoncées. Une simulation avec ces valeurs concorde à chaque $p$. Lorsque $p \to \infty$ avec $n$ fixé, l'erreur tend vers $\frac12$, et lorsque $n \to \infty$, elle tend vers $\Phi(-1)$.

La conclusion~: les variables inutiles nuisent activement. Le bruit dans la direction estimée ajoute $2p/n$ au carré de la longueur, et seule la partie $\Delta^2$ est du signal. Avec 1000 variables et 50 lignes par classe, la direction du classifieur est surtout du bruit. C'est pourquoi la sélection de variables ou la rétraction importe tant lorsque le signal est faible, et un résultat au niveau du hasard n'a pas besoin d'un bogue pour s'expliquer.

\iqwrong{Affirmer que des variables supplémentaires ne peuvent pas nuire, car leurs poids seront simplement proches de zéro. Chaque poids est petit, mais ils sont nombreux, et ensemble ils ont plus de longueur que le signal.}

\iqfollow{À partir de quel $p$ l'erreur atteint-elle 30~\%~?}{$2/\sqrt{4 + p/25} = \Phi^{-1}(0.7) = 0.524$ donne $p \approx 264$. Avec 100 lignes, cela représente environ 260 variables, dont une seule est réelle.}

\iqwhy{Une raison exacte et facile à retenir pour laquelle davantage de variables peut dégrader un modèle. La dérivation n'utilise qu'une moyenne et une variance, une fois le score écrit.}
\end{solution}

\begin{solution}{iq:iv1:machine-learning:45}
La borne provient de la perte exponentielle. Pour $\varepsilon = 0.4$, les deux bornes exigent 226 et 231 tours. Hormis la condition $\varepsilon_t < \frac12$, les données n'interviennent pas.

Soient les étiquettes $\pm 1$ et $F = \sum_t\alpha_th_t$. Les poids partent de $1/n$, et chaque mise à jour multiplie par $e^{-\alpha_ty_ih_t(x_i)}$ puis divise par le facteur de normalisation $Z_t$\index{AdaBoost}. Après $T$ tours,
\[
w_{T + 1}(i) = \frac{e^{-y_iF(x_i)}}{n\prod_tZ_t}.
\]
Les poids somment à 1, donc $\frac1n\sum_ie^{-y_iF(x_i)} = \prod_tZ_t$.

Une erreur vérifie $y_iF(x_i) \le 0$, donc $e^{-y_iF(x_i)} \ge 1$, et l'indicatrice d'erreur est majorée par cette exponentielle. Par conséquent, l'erreur d'entraînement est au plus $\prod_tZ_t$.

Ensuite, $Z_t = (1 - \varepsilon_t)e^{-\alpha_t} + \varepsilon_te^{\alpha_t}$. Le $\alpha_t$ donné minimise cette quantité, avec pour valeur $2\sqrt{\varepsilon_t(1 - \varepsilon_t)}$. Avec $\varepsilon_t = \frac12 - \gamma_t$, elle vaut $\sqrt{1 - 4\gamma_t^2} \le e^{-2\gamma_t^2}$, car $1 - u \le e^{-u}$. Un produit donne la somme dans l'exposant.

Pour $\varepsilon = 0.4$~: $Z = 2\sqrt{0.24} = 0.9798$, et $\ln 0.01/\ln 0.9798$ donne $T = 226$. Pour la seconde borne, $\gamma = 0.1$, donc $T \ge \ln(100)/(2 \times 0.01)$, ce qui donne 231. La borne exponentielle perd très peu, puisque $\gamma$ est petit.

Une exécution avec des souches sur 400 lignes vérifie à chaque tour que $Z_t$ égale sa formule, que la perte exponentielle moyenne égale le produit, et que l'erreur d'entraînement reste inférieure à celui-ci.

La preuve n'exige que chaque apprenant faible batte un tirage à pile ou face sous les poids courants, et aucune hypothèse n'est faite sur la loi des données. Le piège est qu'un tel apprenant peut ne pas exister aux tours suivants, et $\gamma_t$ diminue en général. Par ailleurs, elle ne dit rien de l'erreur de test, qui continue souvent de baisser après que l'erreur d'entraînement a atteint zéro, un phénomène expliqué par les marges.

\iqwrong{Supposer que la borne garantit la généralisation. C'est un énoncé qui ne porte que sur les lignes d'entraînement. Sur des étiquettes financières bruitées, faire tomber l'erreur d'entraînement à zéro est précisément ce qu'il faut éviter.}

\iqfollow{L'erreur d'entraînement ne peut prendre que des valeurs par paliers de $1/n$. Qu'est-ce que cela implique~?}{Dès que la borne passe sous $1/n$, l'erreur d'entraînement doit être exactement nulle. Avec $n = 400$ et $\varepsilon = 0.4$, cela demande au plus 294 tours.}

\iqwhy{C'est un bel exemple de preuve à la force surprenante, et chaque étape est élémentaire. Il teste si le candidat sait passer de la mise à jour de l'algorithme à la borne.}
\end{solution}
